\documentclass[11pt,a4paper]{article}
\input{../../sablony/pripravy/priprava.sty}
\usepackage{tikz}
\nastavPripravu{1. ročník}{Souhrnné opakování}{130}{Logické stopy a množiny}
\begin{document}
\blok{Úvod a připomenutí}
Opakujeme rozpoznání výroku, implikaci, negaci, množinové operace a čtení intervalů. Cíl: argumentovat, \emph{proč} tvrzení platí, nikoli pouze zaškrtávat odpověď. Žádné nové definice. Číselnou osu nech žáky vždy dokreslit k zápisu intervalu.
\blok{Společně řešené příklady}
\begin{prikladReseniBox}[title={Příklad 1}]
Rozhodněte, zda platí tvrzení: „Je-li celé číslo dělitelné šesti, je dělitelné třemi.“ A co obrácené tvrzení?
\tcblower
\[
\begin{aligned}
6\mid n&\ \Rightarrow\ 3\mid n\quad\text{(platí)}\\
3\mid n&\ \not\Rightarrow\ 6\mid n\quad\text{(např. }n=3\text{)}
\end{aligned}
\]

Jeden správný protipříklad stačí k vyvrácení obecného tvrzení.
\end{prikladReseniBox}
\begin{prikladReseniBox}[title={Příklad 2}]
Nechť $A=\{-3;-2;-1;0;1;2;3\}$ a $B=\{-2;0;2;4\}$. Určete průnik, sjednocení a rozdíl $A\setminus B$.
\tcblower
\[
\begin{aligned}
A\cap B&=\{-2;0;2\}\\
A\cup B&=\{-3;-2;-1;0;1;2;3;4\}\\
A\setminus B&=\{-3;-1;1;3\}
\end{aligned}
\]

\end{prikladReseniBox}
\begin{prikladReseniBox}[title={Příklad 3}]
Určete průnik a sjednocení intervalů $I=(-2;3\rangle$ a $J=\langle1;5)$.
\tcblower
\[
\begin{aligned}
I\cap J&=\langle1;3\rangle\\
I\cup J&=(-2;5)
\end{aligned}
\]

\end{prikladReseniBox}
\begin{prikladReseniBox}[title={Příklad 4}]
Hledáme celé číslo $n\in\{2;3;4;6;7;8;9\}$. Právě dvě z tvrzení „$n$ je sudé“, „$n>5$“, „$3\mid n$“ mají být pravdivá. Najděte všechna možná čísla.
\tcblower
\[
\begin{aligned}
n=8&:\quad\text{pravda, pravda, nepravda}\\
n=9&:\quad\text{nepravda, pravda, pravda}\\
K&=\{8;9\}
\end{aligned}
\]

\end{prikladReseniBox}
\blok{Samostatné procvičování}
\begin{zadaniReseniBox}[title={Samostatná práce 1}]
Znegujte výrok: „Každé prvočíslo je liché.“
\tcblower
Negace: „Existuje prvočíslo, které není liché.“ Takovým prvočíslem je $2$.
\end{zadaniReseniBox}
\begin{zadaniReseniBox}[title={Samostatná práce 2}]
Určete $(-5;2)\cap\langle-1;4\rangle$.
\tcblower
\[
\begin{aligned}
K&=\langle-1;2)
\end{aligned}
\]

\end{zadaniReseniBox}
\begin{zadaniReseniBox}[title={Samostatná práce 3}]
Vypište $\{n\in\mathbb Z\mid -2\le n<5\}\setminus\{0;3\}$.
\tcblower
\[
\begin{aligned}
K&=\{-2;-1;1;2;4\}
\end{aligned}
\]

\end{zadaniReseniBox}
\blok{Volitelná rezerva a výzva}
\poznamka{Výzva: sestav dvě různé množiny $A,B$, pro něž je $A\cap B=\varnothing$, ale $A\cup B$ obsahuje osm prvků.}
\blok{Domácí úkol}
\begin{zadaniReseniBox}[title={Domácí úkol}]
\begin{enumerate}[label=\alph*),itemsep=2pt]
\item Vyplňte pravdivostní hodnoty $p\wedge\neg q$ pro dvojice $(p,q)$: $(1,1),(1,0),(0,1),(0,0)$.
\item Která z čísel $\sqrt2,\ -3{,}5,\ 0,\ 5$ jsou racionální?
\item Rozhodněte o platnosti: „Je-li celé číslo dělitelné $4$, je sudé“, a o platnosti obráceného tvrzení.
\item Vypočtěte $\langle-3;2)\cap(0;4\rangle$ a $\langle-3;2)\cup(0;4\rangle$.
\item Vypište všechny prvky $\{x\in\mathbb Z\mid -3\le x<6\}\setminus\{0;4\}$.
\item Najděte všechna celá čísla splňující $0<x<4$ a $x^2-4x<0$.
\end{enumerate}
\tcblower
\begin{enumerate}[label=\alph*),itemsep=2pt]
\item
\[
\begin{aligned}
p\wedge\neg q&:\quad 0,1,0,0
\end{aligned}
\]
\item
\[
\begin{aligned}
\mathbb Q\text{ obsahuje }&-3{,}5,\ 0,\ 5\\
\sqrt2&\notin\mathbb Q
\end{aligned}
\]
\item
\[
\begin{aligned}
4\mid n&\Rightarrow 2\mid n\quad\text{platí}\\
2\mid n&\Rightarrow 4\mid n\quad\text{neplatí (}n=2\text{)}
\end{aligned}
\]
\item
\[
\begin{aligned}
I\cap J&=(0;2)\\
I\cup J&=\langle-3;4\rangle
\end{aligned}
\]
\item
\[
\begin{aligned}
K&=\{-3;-2;-1;1;2;3;5\}
\end{aligned}
\]
\item
\[
\begin{aligned}
x(x-4)&<0\ \Rightarrow\ x\in(0;4)\\
K&=\{1;2;3\}\quad(x\in\mathbb Z)
\end{aligned}
\]
\end{enumerate}
\end{zadaniReseniBox}
\poznamka{Čas: 5 minut zahřátí, 18 minut společný rozbor, 15 minut samostatně, 7 minut obhajoba postupů. Protipříklad je plnohodnotné matematické zdůvodnění.}
\end{document}
